\documentclass[10pt]{beamer}
\usetheme{Madrid}
\usepackage[T1]{fontenc}
\usepackage[utf8]{inputenc}
\usepackage[brazil]{babel}
\usepackage{lmodern}
\usepackage{hyperref}
\usepackage{listings}
\usepackage{xcolor}
\setbeamertemplate{navigation symbols}{}

\lstset{
  language=C++,
  basicstyle=\ttfamily\small,
  keywordstyle=\color{blue},
  commentstyle=\color{gray},
  stringstyle=\color{red},
  breaklines=true,
  showstringspaces=false,
}

\title{Compra de Componentes e Datasheets}
\subtitle{Curso de Microcontroladores}
\author{}
\date{}

\begin{document}

\begin{frame}
  \titlepage
\end{frame}

\begin{frame}{Por que essa aula é importante?}
  \begin{itemize}
    \item Antes de começar a programar, você precisa entender os componentes que vai usar.
    \item Muitos problemas em projetos eletrônicos vêm da \textbf{escolha errada do componente}.
    \item Nem todo sensor GPS é igual. Nem todo motor funciona da mesma forma.
    \item Saber ler um datasheet e entender especificações economiza tempo e dinheiro.
    \item Aprenderemos a identificar um bom fornecedor, fazer perguntas corretas e evitar armadilhas comuns.
  \end{itemize}
\end{frame}

\begin{frame}{Onde comprar componentes eletrônicos?}
  \begin{itemize}
    \item \textbf{Nacionais}: Eletrogate, Robocore, Filipeflop, UFPA Store.
    \item \textbf{Internacionais}: AliExpress, Ebay, Amazon (mais barato, mas demora).
    \item \textbf{Diretos}: Distribuição oficial (mais caro, mas garantido).
    \item Considere: preço, tempo de entrega, política de devolução, reputação do vendedor.
    \item Não compre apenas pelo preço mais baixo! Um componente ruim pode destruir seu projeto.
  \end{itemize}
\end{frame}

\begin{frame}{O que é um Datasheet?}
  \begin{itemize}
    \item Um \textbf{datasheet} é o documento técnico oficial do fabricante de um componente.
    \item Contém todas as especificações, diagramas, pinagem, requisitos de tensão, corrente máxima, etc.
    \item Está \textbf{sempre disponível gratuitamente} na internet (basta procurar pelo modelo).
    \item Exemplo: procure por ``MPU-9250 datasheet PDF'' no Google e você encontrará o documento oficial.
    \item Se não encontrar o datasheet, \textbf{desconfie} do componente. Pode ser falsificado ou descontinuado.
  \end{itemize}
\end{frame}

\begin{frame}{Por que ler o Datasheet?}
  \begin{itemize}
    \item \textbf{Pinagem}: qual pino é VCC, GND, SDA, SCL, etc.
    \item \textbf{Tensão de operação}: funciona em 3.3V ou precisa de 5V?
    \item \textbf{Corrente máxima}: quanto de corrente o componente consome?
    \item \textbf{Protocolos}: usa I2C, SPI, UART ou comunicação digital simples?
    \item \textbf{Especificações técnicas}: precisão, faixa de operação, temperatura máxima, etc.
    \item Sem essas informações, você não consegue conectar o componente corretamente.
  \end{itemize}
\end{frame}

\begin{frame}{Exemplo: Sensor de Temperatura BMP280}
  \begin{itemize}
    \item \textbf{Precisão de Altitude}: ~0.13m.
    \item \textbf{Faixa de Pressão}: 300-1100 hPa.
    \item \textbf{Interface}: SPI ou I2C.
    \item \textbf{Tensão de Operação}: 1.7V-3.6V (compatível com 3.3V).
    \item Do datasheet aprendemos que:
    \begin{itemize}
      \item Funciona com microcontroladores de 3.3V como ESP32.
      \item Usa protocolo I2C (fácil de usar).
      \item Não consome muita corrente (ideal para sistemas embarcados).
    \end{itemize}
  \end{itemize}
\end{frame}

\begin{frame}{Exemplo: Sensor de Distância Ultrassônico}
  \begin{itemize}
    \item Datasheet nos diz:
    \begin{itemize}
      \item Faixa: 2cm a 400cm.
      \item Ângulo de detecção: 15°.
      \item Tempo de resposta: ~60ms.
      \item Pino TRIG (entrada, HIGH por 10µs para disparar).
      \item Pino ECHO (saída, recebe duração do pulso).
    \end{itemize}
    \item Sem essas informações, você não saberia como conectar ou programar o sensor.
    \item Com o datasheet, fica claro exatamente o que fazer.
  \end{itemize}
\end{frame}

\begin{frame}{Módulos vs Componentes Brutos}
  \begin{itemize}
    \item \textbf{Componente bruto}: chip isolado (ex: chip MPU-9250 puro).
    \begin{itemize}
      \item Mais barato, mas exige solda e conhecimento avançado.
      \item Requer circuito de suporte (resistores, capacitores).
    \end{itemize}
    \item \textbf{Módulo}: componente + PCB + circuito de suporte (ex: GY-9250).
    \begin{itemize}
      \item Mais caro, mas muito mais fácil de usar.
      \item Vem com fios prontos (geralmente).
      \item Ideal para iniciantes.
    \end{itemize}
    \item \textbf{Dica}: para começar, sempre compre módulos. Quando tiver experiência, considere componentes brutos.
  \end{itemize}
\end{frame}

\begin{frame}{Tensão de Operação: 3.3V vs 5V}
  \begin{itemize}
    \item Microcontroladores como Arduino tradicional: 5V.
    \item ESP32, ESP8266: 3.3V.
    \item \textbf{Cuidado}: conectar um sensor de 3.3V em 5V pode queimá-lo!
    \item \textbf{Solução}: usar divisor de tensão ou nível shifter.
    \item Sempre verifique a tensão no datasheet antes de comprar.
    \item Componentes de 3.3V são mais comuns em sensores modernos.
  \end{itemize}
\end{frame}

\begin{frame}{Corrente: O que é e por que importa?}
  \begin{itemize}
    \item Cada componente consome uma certa quantidade de corrente (em mA ou µA).
    \item Um microcontrolador pode fornecer apenas uma quantidade limitada de corrente por pino (~40mA em Arduino).
    \item Se você conectar um componente que consome 100mA no pino digital, pode danificar o microcontrolador.
    \item \textbf{Solução}: usar componentes que consomem menos corrente, ou usar um transistor/driver.
    \item Exemplo: LED consome ~20mA (cuidado!), sensor ultrassônico consome ~15mA.
  \end{itemize}
\end{frame}

\begin{frame}{Protocolo de Comunicação}
  \begin{itemize}
    \item Diferentes sensores usam diferentes protocolos:
    \begin{itemize}
      \item \textbf{I2C}: 2 fios (SDA, SCL), vários dispositivos no mesmo barramento.
      \item \textbf{SPI}: 4 fios (MOSI, MISO, CLK, CS), mais rápido que I2C.
      \item \textbf{UART}: comunicação serial, 2 fios (TX, RX).
      \item \textbf{Digital simples}: apenas HIGH/LOW (ex: sensor ultrassônico).
    \end{itemize}
    \item O datasheet especifica qual protocolo o componente usa.
    \item Escolher o protocolo errado significa que o componente não funcionará com seu projeto.
  \end{itemize}
\end{frame}

\begin{frame}{Checklist antes de comprar}
  \begin{enumerate}
    \item \textbf{Encontrou o datasheet?} Se não, não compre.
    \item \textbf{Tensão correta?} Compatível com seu microcontrolador?
    \item \textbf{Protocolo compatível?} I2C/SPI/UART que seu microcontrolador suporta?
    \item \textbf{Pinos disponíveis?} Seu microcontrolador tem os pinos necessários?
    \item \textbf{Corrente?} O microcontrolador consegue fornecer a corrente necessária?
    \item \textbf{Precisão adequada?} O sensor tem precisão suficiente para seu projeto?
    \item \textbf{Vendor confiável?} O vendedor tem boas avaliações?
  \end{enumerate}
\end{frame}

\begin{frame}{Especificações técnicas: O que procurar?}
  \begin{itemize}
    \item \textbf{Faixa de operação}: qual intervalo o sensor consegue medir?
    \item \textbf{Resolução}: qual é o menor valor que o sensor consegue diferenciar?
    \item \textbf{Precisão}: qual é o erro máximo esperado?
    \item \textbf{Tempo de resposta}: quanto tempo demora para o sensor responder?
    \item \textbf{Temperatura de operação}: funciona em ambientes frios/quentes?
    \item Exemplo GPS: precisão de 2.5m significa que pode estar até 2.5m longe da posição real.
  \end{itemize}
\end{frame}

\begin{frame}{Comunicação com o Fornecedor}
  \begin{itemize}
    \item Antes de comprar, você pode fazer perguntas!
    \item \textbf{Em AliExpress}: use o chat, pergunte:
    \begin{itemize}
      \item ``Este módulo é compatível com ESP32 (3.3V)?''
      \item ``Qual é a precisão exata do sensor GPS?''
      \item ``Vem com pinos soldados ou preciso soldar?''
    \end{itemize}
    \item \textbf{Respostas de qualidade}: vendedores bons respondem rápido e com detalhes.
    \item Se o vendedor não responde ou responde vago, compre de outro.
  \end{itemize}
\end{frame}

\begin{frame}{Armadilhas comuns na compra}
  \begin{itemize}
    \item \textbf{Clones/falsificados}: componentes copiados que não funcionam bem.
    \begin{itemize}
      \item Solução: compre de vendedores com muitas avaliações positivas.
    \end{itemize}
    \item \textbf{Descrição enganosa}: anúncio promete precisão mas datasheet diz outra coisa.
    \begin{itemize}
      \item Solução: sempre verifique o datasheet.
    \end{itemize}
    \item \textbf{Falta de suporte técnico}: vendedor desaparece quando há problema.
    \begin{itemize}
      \item Solução: compre de plataformas com proteção do comprador.
    \end{itemize}
    \item \textbf{Tempo de entrega subestimado}: demora muito mais que prometido.
    \begin{itemize}
      \item Solução: leia avaliações sobre tempo de entrega.
    \end{itemize}
  \end{itemize}
\end{frame}

\begin{frame}{Quando você recebe o componente}
  \begin{itemize}
    \item \textbf{Verificar a embalagem}: verifique se chegou em bom estado, sem danos.
    \item \textbf{Verificar o modelo}: compare o modelo no componente com o que pediu.
    \item \textbf{Teste rápido}: antes de integrar ao projeto, faça um teste rápido.
    \begin{itemize}
      \item Conecte o sensor, rode um código de teste simples.
      \item Verifique se as leituras fazem sentido.
    \end{itemize}
    \item \textbf{Se não funcionar}: tente contato com o vendedor rapidamente.
    \item Mantenha a embalagem original por alguns dias (até confirmar que funciona).
  \end{itemize}
\end{frame}

\begin{frame}{Exemplo de Compra Bem Feita}
  \begin{itemize}
    \item \textbf{Projeto}: seguidor de linha com ESP32.
    \item \textbf{Componente necessário}: sensor óptico reflexivo QTR-8.
    \item \textbf{Passos}:
    \begin{enumerate}
      \item Procurei pelo datasheet do QTR-8 (encontrei tensão: 3.3V-5V, interface: analógica).
      \item Verifiquei que ESP32 tem portas analógicas suficientes.
      \item Comprei de um vendedor com muitas avaliações positivas.
      \item Perguntei se vinha com pinos soldados (veio).
      \item Ao chegar, testei com um código simples antes de integrar.
    \end{enumerate}
    \item Resultado: componente funcionou perfeitamente na primeira tentativa.
  \end{itemize}
\end{frame}

\begin{frame}{Quantidade: Compre extras!}
  \begin{itemize}
    \item \textbf{Sempre compre pelo menos 2 unidades de cada componente crítico}.
    \begin{itemize}
      \item Um pode chegar com defeito.
      \item Um pode queimar durante os testes.
      \item Você terá um de backup.
    \end{itemize}
    \item Componentes eletrônicos são baratos, especialmente em quantidade.
    \item Perder 1-2 dias esperando por um novo componente é frustrante.
    \item Exemplos: sensores, módulos de comunicação, motores.
  \end{itemize}
\end{frame}

\begin{frame}{Documentação: Salve tudo!}
  \begin{itemize}
    \item Quando receber o componente, salve o datasheet em uma pasta do projeto.
    \item Tire screenshots de informações importantes no anúncio de compra.
    \item Anote o modelo exato, data de compra, fornecedor.
    \item \textbf{Isso será útil se}:
    \begin{itemize}
      \item Você precisar reproduzir o projeto meses depois.
      \item Alguém pedir para você fazer algo similar.
      \item Você tiver um problema e precisar revisar o datasheet.
    \end{itemize}
    \item Tempo gasto documentando = tempo economizado no futuro.
  \end{itemize}
\end{frame}

\begin{frame}{Resumo: Os 5 passos para comprar corretamente}
  \begin{enumerate}
    \item \textbf{Pesquisar}: datasheet, especificações, avaliações do vendedor.
    \item \textbf{Questionar}: pergunte dúvidas ao vendedor antes de comprar.
    \item \textbf{Verificar compatibilidade}: tensão, protocolo, pinos necessários.
    \item \textbf{Testar}: quando chegar, faça testes antes de integrar ao projeto.
    \item \textbf{Documentar}: salve datasheets e informações importantes.
  \end{enumerate}
  
  \begin{itemize}
    \item Seguir esses passos economiza tempo, dinheiro e muita frustração.
    \item Um projeto bem planejado desde a compra dos componentes tem 80\% de chance de funcionar na primeira tentativa.
  \end{itemize}
\end{frame}

\end{document}
